\section*{Appendix: code excerpts}
\begin{frame}[allowframebreaks]{Code p5.1: k\_F, s and the density floor (exchange network) [xcquinox/alec/networks.py:319--324]}
\label{code:p5_1_k_f_s_and_the_density_floor_exchange_net}
\VerbatimInput[fontsize=\tiny]{slides_code/p5_1_k_f_s_and_the_density_floor_exchange_net.txt}
\end{frame}
\begin{frame}[allowframebreaks]{Code p5.2: r\_s, k\_F, s (correlation network) [xcquinox/alec/networks.py:589--597]}
\label{code:p5_2_r_s_k_f_s_correlation_network}
\VerbatimInput[fontsize=\tiny]{slides_code/p5_2_r_s_k_f_s_correlation_network.txt}
\end{frame}
\begin{frame}[allowframebreaks]{Code p5.3: the DFS coordinates x\_0, x\_1, x\_s (correlation network) [xcquinox/alec/networks.py:626--654]}
\label{code:p5_3_the_dfs_coordinates_x_0_x_1_x_s_correlat}
\VerbatimInput[fontsize=\tiny]{slides_code/p5_3_the_dfs_coordinates_x_0_x_1_x_s_correlat.txt}
\end{frame}
\begin{frame}[allowframebreaks]{Code p5.4: the exchange coordinate x\_s [xcquinox/alec/networks.py:341--352]}
\label{code:p5_4_the_exchange_coordinate_x_s}
\VerbatimInput[fontsize=\tiny]{slides_code/p5_4_the_exchange_coordinate_x_s.txt}
\end{frame}
\begin{frame}[allowframebreaks]{Code p5.5: the log transform and the 1e-5 constant [xcquinox/alec/networks.py:27--34]}
\label{code:p5_5_the_log_transform_and_the_1e_5_constant}
\VerbatimInput[fontsize=\tiny]{slides_code/p5_5_the_log_transform_and_the_1e_5_constant.txt}
\end{frame}
\begin{frame}[allowframebreaks]{Code p5.6: the meta-GGA coordinate ln((alpha + 1)/2) [xcquinox/alec/networks.py:48--80]}
\label{code:p5_6_the_meta_gga_coordinate_ln_alpha_1_2}
\VerbatimInput[fontsize=\tiny]{slides_code/p5_6_the_meta_gga_coordinate_ln_alpha_1_2.txt}
\end{frame}
\begin{frame}[allowframebreaks]{Code p5.7: the exchange network on the doubled spin channel [xcquinox/alec/oneshot.py:465--537]}
\label{code:p5_7_the_exchange_network_on_the_doubled_spin}
\VerbatimInput[fontsize=\tiny]{slides_code/p5_7_the_exchange_network_on_the_doubled_spin.txt}
\end{frame}
\begin{frame}[allowframebreaks]{Code p6.1: tau from the live density matrix [xcquinox/alec/metagga.py:115--139]}
\label{code:p6_1_tau_from_the_live_density_matrix}
\VerbatimInput[fontsize=\tiny]{slides_code/p6_1_tau_from_the_live_density_matrix.txt}
\end{frame}
\begin{frame}[allowframebreaks]{Code p6.2: the smooth positive part [xcquinox/alec/metagga.py:142--152]}
\label{code:p6_2_the_smooth_positive_part}
\VerbatimInput[fontsize=\tiny]{slides_code/p6_2_the_smooth_positive_part.txt}
\end{frame}
\begin{frame}[allowframebreaks]{Code p6.3: alpha = (tau - tau\_W) / tau\_unif [xcquinox/alec/metagga.py:163--274]}
\label{code:p6_3_alpha_tau_tau_w_tau_unif}
\VerbatimInput[fontsize=\tiny]{slides_code/p6_3_alpha_tau_tau_w_tau_unif.txt}
\end{frame}
\begin{frame}[allowframebreaks]{Code p6.4: the alpha descriptor of the network input [xcquinox/alec/descriptors.py:432--472]}
\label{code:p6_4_the_alpha_descriptor_of_the_network_inpu}
\VerbatimInput[fontsize=\tiny]{slides_code/p6_4_the_alpha_descriptor_of_the_network_inpu.txt}
\end{frame}
\begin{frame}[allowframebreaks]{Code p6.5: the cusp pair x\_4, x\_5 [xcquinox/alec/descriptors.py:225--258]}
\label{code:p6_5_the_cusp_pair_x_4_x_5}
\VerbatimInput[fontsize=\tiny]{slides_code/p6_5_the_cusp_pair_x_4_x_5.txt}
\end{frame}
\begin{frame}[allowframebreaks]{Code p7.1: the MLP, GELU, the attention block after the first hidden layer, the gate and the output [xcquinox/alec/networks.py:366--405]}
\label{code:p7_1_the_mlp_gelu_the_attention_block_after_t}
\VerbatimInput[fontsize=\tiny]{slides_code/p7_1_the_mlp_gelu_the_attention_block_after_t.txt}
\end{frame}
\begin{frame}[allowframebreaks]{Code p7.2: the network construction and the zero-initialized last layer [xcquinox/alec/networks.py:283--300]}
\label{code:p7_2_the_network_construction_and_the_zero_in}
\VerbatimInput[fontsize=\tiny]{slides_code/p7_2_the_network_construction_and_the_zero_in.txt}
\end{frame}
\begin{frame}[allowframebreaks]{Code p7.3: the self-attention block [xcquinox/net.py:50--136]}
\label{code:p7_3_the_self_attention_block}
\VerbatimInput[fontsize=\tiny]{slides_code/p7_3_the_self_attention_block.txt}
\end{frame}
\begin{frame}[allowframebreaks]{Code p8.1: the bounded map L\_lambda [xcquinox/alec/networks.py:129--169]}
\label{code:p8_1_the_bounded_map_l_lambda}
\VerbatimInput[fontsize=\tiny]{slides_code/p8_1_the_bounded_map_l_lambda.txt}
\end{frame}
\begin{frame}[allowframebreaks]{Code p8.2: the exchange and correlation energy densities (LDA and PW92 factors) [xcquinox/alec/models.py:175--215]}
\label{code:p8_2_the_exchange_and_correlation_energy_dens}
\VerbatimInput[fontsize=\tiny]{slides_code/p8_2_the_exchange_and_correlation_energy_dens.txt}
\end{frame}
\begin{frame}[allowframebreaks]{Code p8.3: the spin-scaled exchange energy (Oliver and Perdew) [xcquinox/alec/oneshot.py:465--537]}
\label{code:p8_3_the_spin_scaled_exchange_energy_oliver_a}
\VerbatimInput[fontsize=\tiny]{slides_code/p8_3_the_spin_scaled_exchange_energy_oliver_a.txt}
\end{frame}
\begin{frame}[allowframebreaks]{Code p9.1: the architecture record [xcquinox/alec/config.py:103--159]}
\label{code:p9_1_the_architecture_record}
\VerbatimInput[fontsize=\tiny]{slides_code/p9_1_the_architecture_record.txt}
\end{frame}
\begin{frame}[allowframebreaks]{Code p9.2: the registry entries shallow, shallow\_attn, medium and medium\_attn (shown deep\_2x8, deep\_attn\_2x8, deep\_3x16, deep\_attn\_3x16) [xcquinox/alec/config.py:512--515]}
\label{code:p9_2_the_registry_entries_shallow_shallow_att}
\VerbatimInput[fontsize=\tiny]{slides_code/p9_2_the_registry_entries_shallow_shallow_att.txt}
\end{frame}
\begin{frame}[allowframebreaks]{Code p9.3: the registry entries deep\_3x16, deep\_attn\_3x16, deep\_cusp\_3x16 (shown deep0\_*) [xcquinox/alec/config.py:579--592]}
\label{code:p9_3_the_registry_entries_deep_3x16_deep_attn}
\VerbatimInput[fontsize=\tiny]{slides_code/p9_3_the_registry_entries_deep_3x16_deep_attn.txt}
\end{frame}
\begin{frame}[allowframebreaks]{Code p9.4: the registry entry deep\_cusp\_mgga\_3x16 [xcquinox/alec/config.py:691--695]}
\label{code:p9_4_the_registry_entry_deep_cusp_mgga_3x16}
\VerbatimInput[fontsize=\tiny]{slides_code/p9_4_the_registry_entry_deep_cusp_mgga_3x16.txt}
\end{frame}
\begin{frame}[allowframebreaks]{Code p9.5: from\_spec: the fields an entry sets [xcquinox/alec/config.py:385--470]}
\label{code:p9_5_from_spec_the_fields_an_entry_sets}
\VerbatimInput[fontsize=\tiny]{slides_code/p9_5_from_spec_the_fields_an_entry_sets.txt}
\end{frame}
\begin{frame}[allowframebreaks]{Code p10.1: the 21 atomization points (names, spins, charges, references) [xcquinox/alec/dfs\_pool.py:106--225]}
\label{code:p10_1_the_21_atomization_points_names_spins_ch}
\VerbatimInput[fontsize=\tiny]{slides_code/p10_1_the_21_atomization_points_names_spins_ch.txt}
\end{frame}
\begin{frame}[allowframebreaks]{Code p10.2: the three BH76 reaction points [xcquinox/alec/dfs\_pool.py:308--397]}
\label{code:p10_2_the_three_bh76_reaction_points}
\VerbatimInput[fontsize=\tiny]{slides_code/p10_2_the_three_bh76_reaction_points.txt}
\end{frame}
\begin{frame}[allowframebreaks]{Code p10.3: the two ionization points [xcquinox/alec/dfs\_pool.py:427--466]}
\label{code:p10_3_the_two_ionization_points}
\VerbatimInput[fontsize=\tiny]{slides_code/p10_3_the_two_ionization_points.txt}
\end{frame}
\begin{frame}[allowframebreaks]{Code p10.4: the H and Li anchors [xcquinox/alec/dfs\_pool.py:472--477]}
\label{code:p10_4_the_h_and_li_anchors}
\VerbatimInput[fontsize=\tiny]{slides_code/p10_4_the_h_and_li_anchors.txt}
\end{frame}
\begin{frame}[allowframebreaks]{Code p10.5: the pool assembly [xcquinox/alec/dfs\_pool.py:547--639]}
\label{code:p10_5_the_pool_assembly}
\VerbatimInput[fontsize=\tiny]{slides_code/p10_5_the_pool_assembly.txt}
\end{frame}
\begin{frame}[allowframebreaks]{Code p10.6: the 26 points and their species [xcquinox/alec/training\_points.py:347--436]}
\label{code:p10_6_the_26_points_and_their_species}
\VerbatimInput[fontsize=\tiny]{slides_code/p10_6_the_26_points_and_their_species.txt}
\end{frame}
\begin{frame}[allowframebreaks]{Code p11.1: the descriptor triple (rho\^{}(1/3), s, alpha) and the clip to [0, 100] [xcquinox/alec/subset\_selection.py:81--119]}
\label{code:p11_1_the_descriptor_triple_rho_1_3_s_alpha_an}
\VerbatimInput[fontsize=\tiny]{slides_code/p11_1_the_descriptor_triple_rho_1_3_s_alpha_an.txt}
\end{frame}
\begin{frame}[allowframebreaks]{Code p11.2: one PBE SCF per species at def2-svp, grid level 1 [xcquinox/alec/subset\_selection.py:375--406]}
\label{code:p11_2_one_pbe_scf_per_species_at_def2_svp_grid}
\VerbatimInput[fontsize=\tiny]{slides_code/p11_2_one_pbe_scf_per_species_at_def2_svp_grid.txt}
\end{frame}
\begin{frame}[allowframebreaks]{Code p11.3: a point's sample is the concatenation over its species [xcquinox/alec/subset\_selection.py:409--445]}
\label{code:p11_3_a_point_s_sample_is_the_concatenation_ov}
\VerbatimInput[fontsize=\tiny]{slides_code/p11_3_a_point_s_sample_is_the_concatenation_ov.txt}
\end{frame}
\begin{frame}[allowframebreaks]{Code p11.4: the reference histograms: 200 bins between the 0.1 and 99.9 percentiles [xcquinox/alec/subset\_selection.py:448--469]}
\label{code:p11_4_the_reference_histograms_200_bins_betwee}
\VerbatimInput[fontsize=\tiny]{slides_code/p11_4_the_reference_histograms_200_bins_betwee.txt}
\end{frame}
\begin{frame}[allowframebreaks]{Code p11.5: binning with the shared edges and the mass function [xcquinox/alec/subset\_selection.py:175--193]}
\label{code:p11_5_binning_with_the_shared_edges_and_the_ma}
\VerbatimInput[fontsize=\tiny]{slides_code/p11_5_binning_with_the_shared_edges_and_the_ma.txt}
\end{frame}
\begin{frame}[allowframebreaks]{Code p11.6: binning with the shared edges [xcquinox/alec/subset\_selection.py:248--259]}
\label{code:p11_6_binning_with_the_shared_edges}
\VerbatimInput[fontsize=\tiny]{slides_code/p11_6_binning_with_the_shared_edges.txt}
\end{frame}
\begin{frame}[allowframebreaks]{Code p12.1: the probability floor 1e-12 [xcquinox/alec/subset\_selection.py:54--54]}
\label{code:p12_1_the_probability_floor_1e_12}
\VerbatimInput[fontsize=\tiny]{slides_code/p12_1_the_probability_floor_1e_12.txt}
\end{frame}
\begin{frame}[allowframebreaks]{Code p12.2: the Kullback-Leibler term [xcquinox/alec/subset\_selection.py:196--210]}
\label{code:p12_2_the_kullback_leibler_term}
\VerbatimInput[fontsize=\tiny]{slides_code/p12_2_the_kullback_leibler_term.txt}
\end{frame}
\begin{frame}[allowframebreaks]{Code p12.3: the Jensen-Shannon divergence over the three marginals [xcquinox/alec/subset\_selection.py:213--245]}
\label{code:p12_3_the_jensen_shannon_divergence_over_the_t}
\VerbatimInput[fontsize=\tiny]{slides_code/p12_3_the_jensen_shannon_divergence_over_the_t.txt}
\end{frame}
\begin{frame}[allowframebreaks]{Code p13.1: the exhaustive search over C(26, r) [xcquinox/alec/subset\_selection.py:472--688]}
\label{code:p13_1_the_exhaustive_search_over_c_26_r}
\VerbatimInput[fontsize=\tiny]{slides_code/p13_1_the_exhaustive_search_over_c_26_r.txt}
\end{frame}
\begin{frame}[allowframebreaks]{Code p15.1: the pre-training systems: the DFS inventory and the pool atoms [xcquinox/alec/pretrain\_data\_gen.py:334--363]}
\label{code:p15_1_the_pre_training_systems_the_dfs_invento}
\VerbatimInput[fontsize=\tiny]{slides_code/p15_1_the_pre_training_systems_the_dfs_invento.txt}
\end{frame}
\begin{frame}[allowframebreaks]{Code p15.2: the pool atoms [xcquinox/alec/pretrain\_data\_gen.py:214--240]}
\label{code:p15_2_the_pool_atoms}
\VerbatimInput[fontsize=\tiny]{slides_code/p15_2_the_pool_atoms.txt}
\end{frame}
\begin{frame}[allowframebreaks]{Code p15.3: the exchange rows per spin channel at the doubled density [xcquinox/alec/pretrain\_data\_gen.py:379--479]}
\label{code:p15_3_the_exchange_rows_per_spin_channel_at_th}
\VerbatimInput[fontsize=\tiny]{slides_code/p15_3_the_exchange_rows_per_spin_channel_at_th.txt}
\end{frame}
\begin{frame}[allowframebreaks]{Code p15.4: the pointwise targets F\_parent - 1 [xcquinox/alec/pretrain\_data\_gen.py:981--1044]}
\label{code:p15_4_the_pointwise_targets_f_parent_1}
\VerbatimInput[fontsize=\tiny]{slides_code/p15_4_the_pointwise_targets_f_parent_1.txt}
\end{frame}
\begin{frame}[allowframebreaks]{Code p15.5: the synthetic (r\_s, s, alpha) mesh at 30 percent of the weight [xcquinox/alec/pretrain\_data\_gen.py:1068--1140]}
\label{code:p15_5_the_synthetic_r_s_s_alpha_mesh_at_30_per}
\VerbatimInput[fontsize=\tiny]{slides_code/p15_5_the_synthetic_r_s_s_alpha_mesh_at_30_per.txt}
\end{frame}
\begin{frame}[allowframebreaks]{Code p15.6: the integration weights abs(n eps\_LDA) w\_grid [xcquinox/alec/pretrain.py:104--167]}
\label{code:p15_6_the_integration_weights_abs_n_eps_lda_w}
\VerbatimInput[fontsize=\tiny]{slides_code/p15_6_the_integration_weights_abs_n_eps_lda_w.txt}
\end{frame}
\begin{frame}[allowframebreaks]{Code p15.7: the objective: pointwise term plus energy term [xcquinox/alec/pretrain.py:256--341]}
\label{code:p15_7_the_objective_pointwise_term_plus_energy}
\VerbatimInput[fontsize=\tiny]{slides_code/p15_7_the_objective_pointwise_term_plus_energy.txt}
\end{frame}
\begin{frame}[allowframebreaks]{Code p15.8: the learning-rate schedule [xcquinox/alec/pretrain.py:1133--1208]}
\label{code:p15_8_the_learning_rate_schedule}
\VerbatimInput[fontsize=\tiny]{slides_code/p15_8_the_learning_rate_schedule.txt}
\end{frame}
\begin{frame}[allowframebreaks]{Code p15.9: Adam with the global-norm clip [xcquinox/alec/pretrain.py:1211--1236]}
\label{code:p15_9_adam_with_the_global_norm_clip}
\VerbatimInput[fontsize=\tiny]{slides_code/p15_9_adam_with_the_global_norm_clip.txt}
\end{frame}
\begin{frame}[allowframebreaks]{Code p15.10: the loop: validation every validate\_every steps, patience, the best model kept [xcquinox/alec/pretrain.py:740--826]}
\label{code:p15_10_the_loop_validation_every_validate_every}
\VerbatimInput[fontsize=\tiny]{slides_code/p15_10_the_loop_validation_every_validate_every.txt}
\end{frame}
\begin{frame}[allowframebreaks]{Code p16.1: dE\_xc per system in mHa [xcquinox/alec/cluster/fidelity.py:1171--1189]}
\label{code:p16_1_de_xc_per_system_in_mha}
\VerbatimInput[fontsize=\tiny]{slides_code/p16_1_de_xc_per_system_in_mha.txt}
\end{frame}
\begin{frame}[allowframebreaks]{Code p16.2: dAE = dE\_xc(mol) - sum of the atoms' dE\_xc [xcquinox/alec/cluster/fidelity.py:1411--1460]}
\label{code:p16_2_dae_de_xc_mol_sum_of_the_atoms_de_xc}
\VerbatimInput[fontsize=\tiny]{slides_code/p16_2_dae_de_xc_mol_sum_of_the_atoms_de_xc.txt}
\end{frame}
\begin{frame}[allowframebreaks]{Code p16.3: the atomization gate: mean and the max backstop [xcquinox/alec/cluster/fidelity.py:1616--1687]}
\label{code:p16_3_the_atomization_gate_mean_and_the_max_ba}
\VerbatimInput[fontsize=\tiny]{slides_code/p16_3_the_atomization_gate_mean_and_the_max_ba.txt}
\end{frame}
\begin{frame}[allowframebreaks]{Code p16.4: the tolerances 1.0 mHa, 1.0 and 2.0 kcal/mol [xcquinox/alec/cluster/grid\_config.py:290--338]}
\label{code:p16_4_the_tolerances_1_0_mha_1_0_and_2_0_kcal}
\VerbatimInput[fontsize=\tiny]{slides_code/p16_4_the_tolerances_1_0_mha_1_0_and_2_0_kcal.txt}
\end{frame}
\begin{frame}[allowframebreaks]{Code p16.5: the gate that holds the training array [xcquinox/alec/cluster/fidelity.py:292--320]}
\label{code:p16_5_the_gate_that_holds_the_training_array}
\VerbatimInput[fontsize=\tiny]{slides_code/p16_5_the_gate_that_holds_the_training_array.txt}
\end{frame}
\begin{frame}[allowframebreaks]{Code p21.1: the five channels and their assembly [xcquinox/alec/losses.py:1336--1417]}
\label{code:p21_1_the_five_channels_and_their_assembly}
\VerbatimInput[fontsize=\tiny]{slides_code/p21_1_the_five_channels_and_their_assembly.txt}
\end{frame}
\begin{frame}[allowframebreaks]{Code p21.2: the fixed channel weights 1, 1, 1, 1, 20 [xcquinox/alec/train.py:1829--1835]}
\label{code:p21_2_the_fixed_channel_weights_1_1_1_1_20}
\VerbatimInput[fontsize=\tiny]{slides_code/p21_2_the_fixed_channel_weights_1_1_1_1_20.txt}
\end{frame}
\begin{frame}[allowframebreaks]{Code p21.3: the reaction residual (BH76 barriers, W4-11 atomizations as reactions) [xcquinox/alec/losses.py:530--578]}
\label{code:p21_3_the_reaction_residual_bh76_barriers_w4_1}
\VerbatimInput[fontsize=\tiny]{slides_code/p21_3_the_reaction_residual_bh76_barriers_w4_1.txt}
\end{frame}
\begin{frame}[allowframebreaks]{Code p21.4: the BH76 channel [xcquinox/alec/losses.py:1282--1310]}
\label{code:p21_4_the_bh76_channel}
\VerbatimInput[fontsize=\tiny]{slides_code/p21_4_the_bh76_channel.txt}
\end{frame}
\begin{frame}[allowframebreaks]{Code p21.5: the ionization channel [xcquinox/alec/losses.py:581--605]}
\label{code:p21_5_the_ionization_channel}
\VerbatimInput[fontsize=\tiny]{slides_code/p21_5_the_ionization_channel.txt}
\end{frame}
\begin{frame}[allowframebreaks]{Code p21.6: the atom anchors: relative squared error at weight 0.01 [xcquinox/alec/losses.py:223--243]}
\label{code:p21_6_the_atom_anchors_relative_squared_error}
\VerbatimInput[fontsize=\tiny]{slides_code/p21_6_the_atom_anchors_relative_squared_error.txt}
\end{frame}
\begin{frame}[allowframebreaks]{Code p21.7: the atomization channel with the network's own atoms [xcquinox/alec/losses.py:246--273]}
\label{code:p21_7_the_atomization_channel_with_the_network}
\VerbatimInput[fontsize=\tiny]{slides_code/p21_7_the_atomization_channel_with_the_network.txt}
\end{frame}
\begin{frame}[allowframebreaks]{Code p21.8: the potential channel: the Frobenius residual over n\_AO\^{}2 [xcquinox/alec/losses.py:408--483]}
\label{code:p21_8_the_potential_channel_the_frobenius_resi}
\VerbatimInput[fontsize=\tiny]{slides_code/p21_8_the_potential_channel_the_frobenius_resi.txt}
\end{frame}
\begin{frame}[allowframebreaks]{Code p21.9: the density channel per electron [xcquinox/alec/losses.py:366--405]}
\label{code:p21_9_the_density_channel_per_electron}
\VerbatimInput[fontsize=\tiny]{slides_code/p21_9_the_density_channel_per_electron.txt}
\end{frame}
\begin{frame}[allowframebreaks]{Code p21.10: the convergence-tail weights (t/(N-1))\^{}2 [xcquinox/alec/oneshot.py:632--654]}
\label{code:p21_10_the_convergence_tail_weights_t_n_1_2}
\VerbatimInput[fontsize=\tiny]{slides_code/p21_10_the_convergence_tail_weights_t_n_1_2.txt}
\end{frame}
\begin{frame}[allowframebreaks]{Code p21.11: one optimizer step per training group per epoch [xcquinox/alec/train.py:2001--2061]}
\label{code:p21_11_one_optimizer_step_per_training_group_pe}
\VerbatimInput[fontsize=\tiny]{slides_code/p21_11_one_optimizer_step_per_training_group_pe.txt}
\end{frame}
\begin{frame}[allowframebreaks]{Code p21.12: the group's scoped loss [xcquinox/alec/train.py:1919--1965]}
\label{code:p21_12_the_group_s_scoped_loss}
\VerbatimInput[fontsize=\tiny]{slides_code/p21_12_the_group_s_scoped_loss.txt}
\end{frame}
\begin{frame}[allowframebreaks]{Code p22.1: the SCF: three cycles, the DFS mixing schedule, the tail loss [hpcjobs/configs/dfs\_step7.dfs6311\_grid3\_v7g1\_size.yaml:131--140]}
\label{code:p22_1_the_scf_three_cycles_the_dfs_mixing_sche}
\VerbatimInput[fontsize=\tiny]{slides_code/p22_1_the_scf_three_cycles_the_dfs_mixing_sche.txt}
\end{frame}
\begin{frame}[allowframebreaks]{Code p22.2: the mixing schedule a\_t = 0.3\^{}t + 0.3 [xcquinox/alec/solver.py:303--358]}
\label{code:p22_2_the_mixing_schedule_a_t_0_3_t_0_3}
\VerbatimInput[fontsize=\tiny]{slides_code/p22_2_the_mixing_schedule_a_t_0_3_t_0_3.txt}
\end{frame}
\begin{frame}[allowframebreaks]{Code p22.3: the PBE seed [xcquinox/alec/solver.py:88--125]}
\label{code:p22_3_the_pbe_seed}
\VerbatimInput[fontsize=\tiny]{slides_code/p22_3_the_pbe_seed.txt}
\end{frame}
\begin{frame}[allowframebreaks]{Code p22.4: the SCF cycle: the mixed density scored, gradients through every cycle [xcquinox/alec/solver\_manual.py:404--463]}
\label{code:p22_4_the_scf_cycle_the_mixed_density_scored_g}
\VerbatimInput[fontsize=\tiny]{slides_code/p22_4_the_scf_cycle_the_mixed_density_scored_g.txt}
\end{frame}
\begin{frame}[allowframebreaks]{Code p22.5: the cycles as a scan (gradients through every cycle) [xcquinox/alec/solver\_manual.py:263--289]}
\label{code:p22_5_the_cycles_as_a_scan_gradients_through_e}
\VerbatimInput[fontsize=\tiny]{slides_code/p22_5_the_cycles_as_a_scan_gradients_through_e.txt}
\end{frame}
\begin{frame}[allowframebreaks]{Code p22.6: AdamW, the linear decay over the second half, weight decay, the clip [xcquinox/alec/train.py:112--202]}
\label{code:p22_6_adamw_the_linear_decay_over_the_second_h}
\VerbatimInput[fontsize=\tiny]{slides_code/p22_6_adamw_the_linear_decay_over_the_second_h.txt}
\end{frame}
\begin{frame}[allowframebreaks]{Code p23.1: the hyperparameters of the run (200 epochs, 1e-3 to 1e-5, weight decay 1e-4, clip 1.0, validation every 25) [hpcjobs/configs/dfs\_step7.dfs6311\_grid3\_v7g1\_size.yaml:153--183]}
\label{code:p23_1_the_hyperparameters_of_the_run_200_epoch}
\VerbatimInput[fontsize=\tiny]{slides_code/p23_1_the_hyperparameters_of_the_run_200_epoch.txt}
\end{frame}
\begin{frame}[allowframebreaks]{Code p23.2: the validation slice [xcquinox/alec/train.py:745--817]}
\label{code:p23_2_the_validation_slice}
\VerbatimInput[fontsize=\tiny]{slides_code/p23_2_the_validation_slice.txt}
\end{frame}
\begin{frame}[allowframebreaks]{Code p23.3: the validation reaction-energy MAE [xcquinox/alec/train.py:820--862]}
\label{code:p23_3_the_validation_reaction_energy_mae}
\VerbatimInput[fontsize=\tiny]{slides_code/p23_3_the_validation_reaction_energy_mae.txt}
\end{frame}
\begin{frame}[allowframebreaks]{Code p23.4: the validation-best tracker [xcquinox/alec/train.py:682--742]}
\label{code:p23_4_the_validation_best_tracker}
\VerbatimInput[fontsize=\tiny]{slides_code/p23_4_the_validation_best_tracker.txt}
\end{frame}
\begin{frame}[allowframebreaks]{Code p23.5: the validation-best checkpoint among the three saved [xcquinox/alec/train.py:1017--1052]}
\label{code:p23_5_the_validation_best_checkpoint_among_the}
\VerbatimInput[fontsize=\tiny]{slides_code/p23_5_the_validation_best_checkpoint_among_the.txt}
\end{frame}
\begin{frame}[allowframebreaks]{Code p27.1: the full pools (76 BH76, 140 W4-11) before the validation split [xcquinox/alec/full\_benchmark\_pools.py:516--543]}
\label{code:p27_1_the_full_pools_76_bh76_140_w4_11_before}
\VerbatimInput[fontsize=\tiny]{slides_code/p27_1_the_full_pools_76_bh76_140_w4_11_before.txt}
\end{frame}
\begin{frame}[allowframebreaks]{Code p27.2: the validation slice, split by reaction identity [xcquinox/alec/eval\_holdout.py:204--245]}
\label{code:p27_2_the_validation_slice_split_by_reaction_i}
\VerbatimInput[fontsize=\tiny]{slides_code/p27_2_the_validation_slice_split_by_reaction_i.txt}
\end{frame}
\begin{frame}[allowframebreaks]{Code p27.3: the strict held-out filter [xcquinox/alec/eval\_holdout.py:173--201]}
\label{code:p27_3_the_strict_held_out_filter}
\VerbatimInput[fontsize=\tiny]{slides_code/p27_3_the_strict_held_out_filter.txt}
\end{frame}
\begin{frame}[allowframebreaks]{Code p27.4: the cell's trained reactions removed [xcquinox/alec/eval\_holdout.py:283--338]}
\label{code:p27_4_the_cell_s_trained_reactions_removed}
\VerbatimInput[fontsize=\tiny]{slides_code/p27_4_the_cell_s_trained_reactions_removed.txt}
\end{frame}
\begin{frame}[allowframebreaks]{Code p27.5: the energy legs average one term per reaction identity [xcquinox/alec/eval\_holdout.py:402--442]}
\label{code:p27_5_the_energy_legs_average_one_term_per_rea}
\VerbatimInput[fontsize=\tiny]{slides_code/p27_5_the_energy_legs_average_one_term_per_rea.txt}
\end{frame}
\begin{frame}[allowframebreaks]{Code p27.6: eps\_n per electron [xcquinox/alec/evaluation.py:196--218]}
\label{code:p27_6_eps_n_per_electron}
\VerbatimInput[fontsize=\tiny]{slides_code/p27_6_eps_n_per_electron.txt}
\end{frame}
\begin{frame}[allowframebreaks]{Code p27.7: the grid RMSE [xcquinox/alec/evaluation.py:293--320]}
\label{code:p27_7_the_grid_rmse}
\VerbatimInput[fontsize=\tiny]{slides_code/p27_7_the_grid_rmse.txt}
\end{frame}
\begin{frame}[allowframebreaks]{Code p27.8: WTMAD-2 with the GMTKN55 scale [notebooks/analysis/make\_ablation\_arch\_figure.py:3824--3857]}
\label{code:p27_8_wtmad_2_with_the_gmtkn55_scale}
\VerbatimInput[fontsize=\tiny]{slides_code/p27_8_wtmad_2_with_the_gmtkn55_scale.txt}
\end{frame}
\begin{frame}[allowframebreaks]{Code p27.9: the harmonic mean [notebooks/analysis/make\_ablation\_arch\_figure.py:4835--4841]}
\label{code:p27_9_the_harmonic_mean}
\VerbatimInput[fontsize=\tiny]{slides_code/p27_9_the_harmonic_mean.txt}
\end{frame}
\begin{frame}[allowframebreaks]{Code p27.10: the self-calibrated gamma = E\_PBE / D\_PBE [notebooks/analysis/make\_ablation\_arch\_figure.py:4868--4943]}
\label{code:p27_10_the_self_calibrated_gamma_e_pbe_d_pbe}
\VerbatimInput[fontsize=\tiny]{slides_code/p27_10_the_self_calibrated_gamma_e_pbe_d_pbe.txt}
\end{frame}
\begin{frame}[allowframebreaks]{Code p27.11: the DFS gamma 1084.87 [notebooks/analysis/make\_ablation\_arch\_figure.py:4952--4975]}
\label{code:p27_11_the_dfs_gamma_1084_87}
\VerbatimInput[fontsize=\tiny]{slides_code/p27_11_the_dfs_gamma_1084_87.txt}
\end{frame}
